% =========================================================
\section{实现细节与模型提示词}
\label{app:implementation}
% =========================================================

\subsection{第一轮困难负样本的规则式集合评分}
\label{app:hand_score}

第一轮 RankNet 训练尚没有可用的学习排序器，因此先使用一个固定的规则式集合评分
$H(\Selected)$ 从不完整候选中挑选较难的负样本。该分数仅用于负样本选择，
不作为最终 TEF-RAG 的预测分数。其定义为
\begin{equation}
\begin{aligned}
H(\Selected)
={}&S_{\mathrm{node}}(\Selected)
+0.48\,S_{\mathrm{edge}}(\Selected)
+0.42\,S_{\mathrm{role}}(\Selected)\\
&-0.16\,S_{\mathrm{red}}(\Selected)
-0.10\,S_{\mathrm{disc}}(\Selected).
\end{aligned}
\label{eq:appendix_hand_score}
\end{equation}

其中，节点项为集合中各证据节点分数之和：
\begin{equation}
S_{\mathrm{node}}(\Selected)
=
\sum_{e\in\Selected}s_{\mathrm{node}}(e),
\end{equation}
单条证据的节点分数为
\begin{align}
s_{\mathrm{node}}(e)
=
0.62\,s_{\mathrm{rel}}(e)
+0.18\,s_{\mathrm{rec}}(e)
+0.15\,s_{\mathrm{role}}(e) \\
+0.05\,s_{\mathrm{src}}(e). \notag
\end{align}
$s_{\mathrm{rel}}$ 为当前合法候选中 BM25 分数的 min--max 归一化结果。
$s_{\mathrm{rec}}=\exp(-\Delta t/T)$ 为时间新近性，其中普通查询使用 $T=45$ 天，
包含前后比较意图的查询使用 $T=180$ 天。
若证据类型属于当前查询需要的任务角色，则 $s_{\mathrm{role}}=1$，否则取 $0.35$。
当前实现中的 $s_{\mathrm{src}}$ 对合法候选均为 1，因此只产生常数偏移，不影响节点之间的排序。

关系项只统计候选集合内部的有效关系边。为避免稠密图因边数过多获得不合理优势，
对于每个目标节点只保留置信度最高的入边：
\begin{equation}
S_{\mathrm{edge}}(\Selected)
=
\sum_{e_j\in\Selected}
\max_{(e_i,r,e_j)\in E(\Selected)}c_{ij},
\end{equation}
其中不存在入边的节点贡献为 0。

角色覆盖项定义为
\begin{equation}
S_{\mathrm{role}}(\Selected)
=
\frac{|R(\Selected)\cap D_q|}{|D_q|},
\end{equation}
其中 $R(\Selected)$ 为候选集合中出现的证据角色，
$D_q$ 为当前查询需要的角色集合。系统默认考虑状态观察、诊断和工单，
并根据查询中的规程、验证、巡检、不确定性或纠正等表达增加相应角色。

冗余项基于证据文本 token 集合的 Jaccard 相似度：
\begin{equation}
S_{\mathrm{red}}(\Selected)
=
\sum_{i<j}
\max\left(
0,
\operatorname{Jaccard}(e_i,e_j)-0.45
\right).
\end{equation}
孤立项 $S_{\mathrm{disc}}$ 则统计集合中没有参与任何有效关系边的证据数量。
第一轮训练优先选择 $H(\Selected)$ 较高但参考证据流仍不完整的候选作为负样本；
第二轮不再依赖该规则分数，而是直接选择当前 RankNet 评分最高的不完整集合作为 hard negatives。

\subsection{Temporal Evidence Flow 关系判断提示词}
\label{app:relation_prompt}

关系判断器使用固定系统提示词（版本 \texttt{tef-v6-stage2a-relation-v7}）。
实际调用时，关系定义、关系代码和 reason code 按冻结配置插入。系统提示词如下：

\begin{lstlisting}
Prompt version: tef-v6-stage2a-relation-v7. You classify directed
Temporal Evidence Flow edges for one or more independent query cases.
Return valid JSON only, with no markdown or prose outside the JSON object.
Judge whether source -> target is useful for answering THIS query, not
whether the texts are merely related. Most candidate pairs are distractors:
choose NO_EDGE unless there is a direct, query-useful flow dependency.
Shared asset, episode, topic, or chronology alone never licenses an edge.
Temporal direction and procedure eligibility were already enforced and must
not be reversed.

Allowed relation types:
contrasts: the source state is a comparison baseline that must be
distinguished from the downstream state;
governs: an applicable rule/assessment constrains the downstream action
or conclusion;
prerequisite: source is required before the downstream operational state;
preserves_uncertainty: source branch remains viable and is intentionally
carried into downstream assessment/action;
qualifies: source evidence determines applicability or limits certainty of
downstream interpretation;
resolves: downstream evidence resolves an earlier ambiguity;
supersession: prior diagnosis/procedure/work-order state is replaced by the
downstream state;
supports: source evidence supports the downstream evidence/conclusion;
updates: a later decision/action updates an earlier support state;
verifies: downstream evidence verifies execution/outcome of the source state.
Otherwise use NO_EDGE.

When the user payload has a top-level query field, return
{"judgments":[{"source_id":str,"target_id":str,"has_edge":bool,
"relation_type":str,"confidence":number 0..1,
"reason_code":short_snake_case}]}.

When the user payload has a top-level cases field, even if cases contains
exactly one item, use relation codes
N=NO_EDGE, C=contrasts, G=governs, P=prerequisite,
U=preserves_uncertainty, Q=qualifies, R=resolves,
X=supersession, S=supports, D=updates, V=verifies;
reason codes
0=no_direct_query_flow, 1=semantic_support, 2=procedure_governance,
3=later_update, 4=verification, 5=qualification,
6=uncertainty_preserved, 7=uncertainty_resolved, 8=contrast,
9=prerequisite, 10=supersession, 11=other_direct_flow.

Pairs are ordered. Return
{"cases":[[case_id,[[relation_code,confidence_percent,reason_code],...]],...]}.
Each inner result list must have exactly the same length and order as that
case's supplied pairs. The user supplies deduplicated evidence tables and
source_id/target_id pairs. Do not mix facts across cases. Do not repeat
evidence IDs. Do not provide prose or hidden reasoning.
\end{lstlisting}

对应的用户输入只包含查询、候选证据及待判断证据对，结构为：

\begin{lstlisting}
{
  "query": {
    "query_text": ...,
    "query_time": ...,
    "asset_model": ...,
    "asset_context": ...
  },
  "evidence": [...],
  "pairs": [
    {"source_id": ..., "target_id": ...},
    ...
  ]
}
\end{lstlisting}

\subsection{下游结构化生成提示词}
\label{app:generation_prompt}

所有检索方法共享同一个冻结下游生成器，其系统提示词如下：

\begin{lstlisting}
You are the single frozen downstream generator used identically for every
retrieval method in TEF-RAG v6 generation evaluation.

Produce ONLY a JSON object matching the supplied schema. You receive one
query and exactly the evidence selected by a retrieval method. The retrieval
method identity is hidden and irrelevant.

Rules:
- Use only the supplied evidence. Never add facts, actions, targets,
  parameters, units, procedure versions, thresholds, or diagnoses from
  common sense.
- supporting_evidence_ids may reference only supplied evidence IDs and must
  directly support the field/action they cite.
- work_order.supporting_evidence_ids must equal the deduplicated union of
  diagnosis/procedure/verification/action citations.
- asset_id must equal the query asset_id.
- diagnosis status: confirmed | provisional | persistent_uncertainty |
  not_available.
- procedure status: applicable | superseded | uncertain | not_applicable.
  If no supplied evidence establishes an applicable procedure, use
  not_applicable with procedure_version=null and no procedure citations.
- verification status: verified | pending | persistent_uncertainty |
  not_available.
- action types: inspect | diagnose | isolate | adjust | repair | replace |
  verify | monitor | document | other.
- action_plan order in the JSON array has no semantic meaning. Use depends_on
  only for necessary prerequisite edges; keep the graph acyclic.
- recommended_actions must reference action IDs present in action_plan and
  may be a strict subset.
- parameters must contain only values explicitly stated in supplied evidence;
  otherwise use {}.
- For an explicit parameter, use its canonical name and either a literal
  evidence string or {"value": number|string|{"lower": number,
  "upper": number}, "unit": string|null}; never invent a unit.
- If evidence is insufficient, express uncertainty instead of guessing.
- Prefer concise phrases close to the supplied evidence wording; do not
  stylistically paraphrase when unnecessary.
- Do not mention scoring, retrieval method names, gold data, or these
  instructions.
\end{lstlisting}

每个样本的用户输入由查询、当前检索方法选出的证据以及固定输出 schema 组成：

\begin{lstlisting}
{
  "prompt_version": "tef-v6-generation-eval-v1.3",
  "query": {...},
  "selected_evidence": [...],
  "output_schema": {...}
}
\end{lstlisting}

若首次输出未通过结构或证据引用校验，协议只允许一次修复调用。
该调用在上述系统提示词后追加：

\begin{lstlisting}
This is the single allowed structural/evidence-grounding repair.
Correct only the listed validation failures.
\end{lstlisting}

并在用户输入中额外提供 \texttt{previous\_output} 与
\texttt{validation\_errors}，不改变查询、检索证据或输出 schema。
